En un laboratori s'ha fet l'experiment que es mostra a la figura 1. En una cubeta de plàstic transparent s'ha afegit aproximadament un centímetre d'aigua de l'aixeta, i s'han col·locat a banda i banda dues plaques conductores de coure separades a una distància de 20~cm. Les plaques s'han connectat a una font d'alimentació. A sota de la cubeta transparent hi ha un paper quadriculat que permet determinar les posicions (figura 1). A la placa connectada al terminal negatiu de la font d'alimentació s'hi ha connectat el terminal negatiu del voltímetre. El terminal positiu del voltímetre s'ha mogut per diferents punts de la quadrícula per a mesurar el potencial elèctric i el resultat s'indica a la taula de sota.

\figura[0.8]{fig1.png}
\begin{center}
\textsc{Figura 1}. Esquema del muntatge de l'experiment i paper quadriculat de la cubeta
\end{center}

\begin{center}
\begin{tabular}{|c|c|c|c|c|}
\hline
\textit{x} (cm) & $V_\mathrm{I}$ (V) & $V_\mathrm{II}$ (V) & $V_\mathrm{III}$ (V) & $V_\mathrm{mitj\grave{a}}$ (V) \\
\hline
4,00 & 1,4 & 1,5 & 1,4 & \\
\hline
8,00 & 2,8 & 2,9 & 2,8 & \\
\hline
12,00 & 4,2 & 4,3 & 4,4 & \\
\hline
16,00 & 5,7 & 5,7 & 5,8 & \\
\hline
\end{tabular}
\end{center}

\begin{apartat}{1,25}
Empleneu la taula de dalt amb la mitjana aritmètica del potencial elèctric a les posicions $x = 4$, 8, 12 i 16~cm. Dibuixeu les línies equipotencials a $x = 4$, 8, 12 i 16~cm i les línies de camp elèctric en el paper quadriculat de la cubeta (figura 1). Representeu en els eixos de coordenades (figura 2) la mitjana aritmètica del potencial elèctric en funció de $x$. Calculeu el mòdul del camp elèctric a partir de la gràfica.
\end{apartat}

\figura[0.45]{fig2.png}
\begin{center}
\textsc{Figura 2}
\end{center}

\begin{apartat}{1,25}
Col·loquem una càrrega positiva de 3,00~mC al punt $P$ indicat dins la cubeta en la figura 1. Indiqueu quina trajectòria seguirà. Representeu en el paper quadriculat (figura 1) la direcció i el sentit de la força que aplica el camp elèctric sobre aquesta càrrega. Determineu el mòdul de la força. Calculeu el treball que fa el camp elèctric per moure la càrrega des de $x = 8$~cm fins a $x = 0$~cm.
\end{apartat}
